% GENERATED FILE. Edit canonical lessons, then run npm run generate:books.
% canonical-source-hash: fnv1a64:691577773fe2affd
% canonical-lessons: GU-C420-kamad, GU-C420-naliyu, GU-C420-vad, GU-C420-thabhlo, GU-C420-chavu

\chapter{Small Houses and Fences}
\label{ch:gu-small-houses-and-fences}
\hlchaptermodality{420}

\begin{chapteropening}
\textbf{By the end of this chapter:} \emph{I can describe a simple country house.}

Complete the last lesson of chapter 420: I can describe a simple country house.
\end{chapteropening}

\section[\texorpdfstring{\gu{કમાડ} (kamāḍ)}{કમાડ (kamāḍ)}]{\gu{કમાડ} (kamāḍ) — a wooden door; a door leaf}

\label{lesson:GU-C420-kamad}



\begin{quote}
Before the new one: say the Gujarati for a resident, then the Gujarati for to wander aimlessly.
\end{quote}

\subsection*{You'll want to know: \gu{કમાડ}}
\textbf{\gu{કમાડ}} — \emph{kamāḍ} — \textquotedblleft{}a wooden door; a door leaf\textquotedblright{}.

\begin{quote}
\gu{પવનથી} \gu{ઘરનું} \gu{જૂનું} \gu{કમાડ} \gu{ખખડે} \gu{છે}.
\end{quote}

\begin{quote}
\emph{pavanthī gharnũ jūnũ kamāḍ khakhḍe che.}
\end{quote}

\begin{quote}
The old wooden door of the house rattles in the wind.
\end{quote}

A neuter noun.

\begin{cousinweb}
\textbf{\gu{કમાડ}} (\emph{kamāḍ}) is a tadbhava of Sanskrit \emph{kavāṭa}, \textquotedblleft{}the leaf of a door\textquotedblright{}; CDIAL lists the Gujarati word, \textquotedblleft{}a door\textquotedblright{} (CDIAL 2963).
\end{cousinweb}

\begin{grammarlens}[title={What you've built}]
One more word for getting around.
\end{grammarlens}

\subsection*{Your turn}
\begin{itemize}
\raggedright
  \item \emph{Say it:} \emph{kamāḍ}
  \item \emph{Say it:} \emph{kamāḍ}, once more
  \item \emph{Say it:} \emph{pavanthī gharnũ jūnũ kamāḍ khakhḍe che.} — The old wooden door of the house rattles in the wind.
  \item \emph{Recall:} say the Gujarati for a resident, then the Gujarati for to wander aimlessly, then say \emph{kamāḍ} again
\end{itemize}

\begin{tcolorbox}[breakable,colback=teal!4,colframe=teal!35!black,title={Before you move on}]
What does \gu{કમાડ} mean? (\textquotedblleft{}A wooden door; a door leaf\textquotedblright{}.) Say it once more.
\end{tcolorbox}
\section[\texorpdfstring{\gu{નળિયું} (naḷiyũ)}{નળિયું (naḷiyũ)}]{\gu{નળિયું} (naḷiyũ) — a roof tile (a curved clay tile)}

\label{lesson:GU-C420-naliyu}



\begin{quote}
Before the new one: say the Gujarati for to wander aimlessly, then the Gujarati for a wooden door.
\end{quote}

\subsection*{You'll want to know: \gu{નળિયું}}
\textbf{\gu{નળિયું}} — \emph{naḷiyũ} — \textquotedblleft{}a roof tile (a curved clay tile)\textquotedblright{}.

\begin{quote}
\gu{તોફાનમાં} \gu{ઘરનાં} \gu{ઘણાં} \gu{નળિયાં} \gu{તૂટી} \gu{ગયાં}.
\end{quote}

\begin{quote}
\emph{tophānmā̃ gharnā̃ ghaṇā̃ naḷiyā̃ tūṭī gayā̃.}
\end{quote}

\begin{quote}
Many roof tiles of the house broke in the storm.
\end{quote}

A neuter noun.

\begin{cousinweb}
\textbf{\gu{નળિયું}} (\emph{naḷiyũ}) seems to be built on \gu{નળ}, a pipe, which goes back to Sanskrit \emph{naḍa}, \textquotedblleft{}a reed\textquotedblright{} (CDIAL 6936): a curved tile is like half a pipe.
\end{cousinweb}

\begin{grammarlens}[title={What you've built}]
One more word for getting around.
\end{grammarlens}

\subsection*{Your turn}
\begin{itemize}
\raggedright
  \item \emph{Say it:} \emph{naḷiyũ}
  \item \emph{Say it:} \emph{naḷiyũ}, once more
  \item \emph{Say it:} \emph{tophānmā̃ gharnā̃ ghaṇā̃ naḷiyā̃ tūṭī gayā̃.} — Many roof tiles of the house broke in the storm.
  \item \emph{Recall:} say the Gujarati for to wander aimlessly, then the Gujarati for a wooden door, then say \emph{naḷiyũ} again
\end{itemize}

\begin{tcolorbox}[breakable,colback=teal!4,colframe=teal!35!black,title={Before you move on}]
What does \gu{નળિયું} mean? (\textquotedblleft{}A roof tile (a curved clay tile)\textquotedblright{}.) Say it once more.
\end{tcolorbox}
\section[\texorpdfstring{\gu{વાડ} (vāḍ)}{વાડ (vāḍ)}]{\gu{વાડ} (vāḍ) — a fence; a hedge}

\label{lesson:GU-C420-vad}



\begin{quote}
Before the new one: say the Gujarati for a wooden door, then the Gujarati for a roof tile.
\end{quote}

\subsection*{You'll want to know: \gu{વાડ}}
\textbf{\gu{વાડ}} — \emph{vāḍ} — \textquotedblleft{}a fence; a hedge\textquotedblright{}.

\begin{quote}
\gu{બગીચાની} \gu{બહાર} \gu{લીલી} \gu{વાડ} \gu{છે}.
\end{quote}

\begin{quote}
\emph{bagīcānī bahār līlī vāḍ che.}
\end{quote}

\begin{quote}
There is a green hedge outside the garden.
\end{quote}

A feminine noun.

\begin{cousinweb}
\textbf{\gu{વાડ}} (\emph{vāḍ}) is a tadbhava of Sanskrit \emph{vāṭa}, \textquotedblleft{}an enclosure\textquotedblright{}; CDIAL lists the Gujarati word, \textquotedblleft{}a fence\textquotedblright{} (CDIAL 11480).
\end{cousinweb}

\begin{grammarlens}[title={What you've built}]
One more word for getting around.
\end{grammarlens}

\subsection*{Your turn}
\begin{itemize}
\raggedright
  \item \emph{Say it:} \emph{vāḍ}
  \item \emph{Say it:} \emph{vāḍ}, once more
  \item \emph{Say it:} \emph{bagīcānī bahār līlī vāḍ che.} — There is a green hedge outside the garden.
  \item \emph{Recall:} say the Gujarati for a wooden door, then the Gujarati for a roof tile, then say \emph{vāḍ} again
\end{itemize}

\begin{tcolorbox}[breakable,colback=teal!4,colframe=teal!35!black,title={Before you move on}]
What does \gu{વાડ} mean? (\textquotedblleft{}A fence; a hedge\textquotedblright{}.) Say it once more.
\end{tcolorbox}
\section[\texorpdfstring{\gu{થાંભલો} (thā̃bhlo)}{થાંભલો (thā̃bhlo)}]{\gu{થાંભલો} (thā̃bhlo) — a post; a pole}

\label{lesson:GU-C420-thabhlo}



\begin{quote}
Before the new one: say the Gujarati for a roof tile, then the Gujarati for a hedge.
\end{quote}

\subsection*{You'll want to know: \gu{થાંભલો}}
\textbf{\gu{થાંભલો}} — \emph{thā̃bhlo} — \textquotedblleft{}a post; a pole\textquotedblright{}.

\begin{quote}
\gu{રસ્તા} \gu{પરના} \gu{થાંભલા} \gu{પર} \gu{બત્તી} \gu{છે}.
\end{quote}

\begin{quote}
\emph{rastā parnā thā̃bhlā par battī che.}
\end{quote}

\begin{quote}
There is a light on the post by the road.
\end{quote}

A masculine noun.

\begin{cousinweb}
\textbf{\gu{થાંભલો}} (\emph{thā̃bhlo}) is a tadbhava of Sanskrit \emph{stambha}, \textquotedblleft{}a post, a pillar\textquotedblright{}; CDIAL lists the Gujarati word (CDIAL 13682).
\end{cousinweb}

\begin{grammarlens}[title={What you've built}]
One more word for getting around.
\end{grammarlens}

\subsection*{Your turn}
\begin{itemize}
\raggedright
  \item \emph{Say it:} \emph{thā̃bhlo}
  \item \emph{Say it:} \emph{thā̃bhlo}, once more
  \item \emph{Say it:} \emph{rastā parnā thā̃bhlā par battī che.} — There is a light on the post by the road.
  \item \emph{Recall:} say the Gujarati for a roof tile, then the Gujarati for a hedge, then say \emph{thā̃bhlo} again
\end{itemize}

\begin{tcolorbox}[breakable,colback=teal!4,colframe=teal!35!black,title={Before you move on}]
What does \gu{થાંભલો} mean? (\textquotedblleft{}A post; a pole\textquotedblright{}.) Say it once more.
\end{tcolorbox}
\section[\texorpdfstring{\gu{છાવું} (chāvũ)}{છાવું (chāvũ)}]{\gu{છાવું} (chāvũ) — to thatch; to roof (a house)}

\label{lesson:GU-C420-chavu}



\begin{quote}
Before the new one: say the Gujarati for a hedge, then the Gujarati for a post.
\end{quote}

\subsection*{You'll want to know: \gu{છાવું}}
\textbf{\gu{છાવું}} — \emph{chāvũ} — \textquotedblleft{}to thatch; to roof (a house)\textquotedblright{}.

\begin{quote}
\gu{દાદાએ} \gu{ઘાસથી} \gu{ઘર} \gu{છાયું}.
\end{quote}

\begin{quote}
\emph{dādāe ghāsthī ghar chāyũ.}
\end{quote}

\begin{quote}
Grandfather thatched the house with grass.
\end{quote}

Past, with \gu{તેણે} (he or she …): \gu{તેણે} \gu{ઘર} \gu{છાયું}, \emph{teṇe ghar chāyũ}.

\begin{cousinweb}
\textbf{\gu{છાવું}} (\emph{chāvũ}) is a tadbhava of Sanskrit \emph{chādáyati}, \textquotedblleft{}covers, veils\textquotedblright{}; CDIAL lists the Gujarati verb, \textquotedblleft{}to thatch\textquotedblright{} (CDIAL 5018): to cover a house with a roof.
\end{cousinweb}

\begin{grammarlens}[title={What you've built}]
One more word for getting around.
\end{grammarlens}

\subsection*{Your turn}
\begin{itemize}
\raggedright
  \item \emph{Say it:} \emph{chāvũ}
  \item \emph{Say it:} \emph{chāvũ}, once more
  \item \emph{Say it:} \emph{dādāe ghāsthī ghar chāyũ.} — Grandfather thatched the house with grass.
  \item \emph{Recall:} say the Gujarati for a hedge, then the Gujarati for a post, then say \emph{chāvũ} again
\end{itemize}

\begin{tcolorbox}[breakable,colback=teal!4,colframe=teal!35!black,title={Before you move on}]
What does \gu{છાવું} mean? (\textquotedblleft{}To thatch; to roof (a house)\textquotedblright{}.) Say it once more.
\end{tcolorbox}
